\section{Validation}
\label{sec:validation}

\paragraph{Closure.}
Three existing closure families were exercised (1A lgbm 5-iter; thresholds in parentheses):
\begin{itemize}
  \item \textbf{Truth-reweight} (multiply truth and reco weights by the same
        factor): gaussian-$\pT$ residual \SI{0.046}{\percent}, tilt-$\pPar$
        \SI{0.013}{\percent} (thr \SI{1.5}{\percent}) --- injected truth recovered.
  \item \textbf{Hidden-variable} (bias on $\Delta\pT=\pT^{\mathrm{reco}}-\pT^{\mathrm{truth}}$,
        a resolution variable not in the feature set): leakage linear in amplitude,
        $A=0.05/0.10/0.30\to0.57/1.14/3.43\,\%$ (thr \SI{3.0}{\percent});
        the first two amplitudes pass and the deliberately large $A=0.30$ arm exceeds
        the threshold.
  \item \textbf{Alt-model} (train on CV, pseudo-data from an alternative universe):
        MaCCQE \SI{0.068}{\percent}, Flux \SI{0.376}{\percent} (thr \SI{2.0}{\percent}).
\end{itemize}

\paragraph{Bottom-line test.}
Following the unbinned-unfolding guide~\cite{unbinnedguide}, the unfolded result must
recover the data's deviation from the prior with a residual far smaller than the feature
itself (the unfolding cannot manufacture discriminating power). Using the closure runs,
in the bins carrying the injected feature the unfolded result tracks the injected truth to
$\approx10\times$ better than the feature size: 2D (gaussian-bump truth reweight) residual
\SI{1.69}{\percent} vs injected \SI{17.2}{\percent} (ratio $0.098$); 3D (+\SI{30}{\percent}
$\Eavail$ bump, $\Eavail$ projection) \SI{1.84}{\percent} vs \SI{18.1}{\percent} (ratio
$0.102$). Both pass.
(The naive data-vs-prior form---reco vs truth $\chi^2$---instead has truth$>$reco, as
expected for a high-resolution measurement where detector smearing lowers the reco baseline
and a stat-only diagonal omits the unfolding covariance; the proper full-covariance
goodness-of-fit is the $\chi^2$-vs-paper of \S\ref{sec:results-tension} and the generator
comparisons.)

These closure tests validate the estimator but are not used to assign a
non-closure uncertainty. Such an uncertainty would require a predeclared suite
of injection shapes and amplitudes rather than calibration to a single
Gaussian bump.

They also do not test the distinct failure mode in which pseudo-data are
generated with a detector response that disagrees with the response supplied
to the unfold. Appendix~\ref{app:response-mismatch} defines a complementary
record-only $\Eavail$ response-mismatch closure and a more complete
leading-nucleon follow-up. No result from that test is claimed here.

Nor do they exercise the background subtraction: every closure above reweights the
simulated signal and unfolds it without background. End-to-end five-dimensional
pseudo-experiments that add background-simulation events and rebuild the per-bin purity
weights exactly as the measurement does (scalar-5D campaign, development study D1) return
the highest-$W$ column of the $(\Eavail,W)$ projection \SIrange{3.5}{4.1}{\percent} low
and the highest-$\Eavail$, highest-$W$ cell \SI{1.7}{\percent} high at nominal truth,
while the same experiments without background close to \SI{0.22}{\percent}. Set against
the quoted total uncertainty of those cells this is 0.36--1.07 standard deviations. The
bias is a property of the purity footing (\S\ref{app:negweight}), is not in any quoted
covariance, and its size on data is not established.
% Sources: docs/orchestration/OUTCOME-20260925-s5c-purity-background-bias-at-high-W.md,
% docs/orchestration/state/s5c/d1/d1_summary.json, docs/orchestration/state/s5c/d1/bias_vs_adopted.json.

\paragraph{Unbinned background treatment and generator-sized truth changes.}
The same end-to-end experiments were repeated with the unbinned negative-weight
treatment (\texttt{-{}-bkg-mode negweight-refined}): data and signed background
simulation are refined to non-negative weights before the first OmniFold step, and the
background simulation is split into an injection half and an independent subtraction
half. Sixty experiments at nominal truth bring the highest-$W$ column to
\SIrange{0.11}{0.24}{\percent} and the corner cell to \SI{0.24}{\percent}. A set of joint
cells nevertheless remains \SIrange{0.7}{1.1}{\percent} high, up to 14 standard errors,
against \SI{-0.2}{\percent} without background; the mechanism is not established.
Experiments whose truth is reweighted by the GiBUU-to-GENIE $\Eavail$ shape ratio
(0.70--1.33 per $\Eavail$ bin, total preserved) are not followed by the unfolding. The
$(\Eavail,W)$ cells deviate from that truth by a median of \SI{9.7}{\percent}, up to
\SI{74}{\percent} in the highest-$W$ column, and by several percent in the largest cells.
A $q_3$ deformation that leaves the $(\Eavail,W)$ truth unchanged moves those cells by up
to \SI{37}{\percent}. Signal-only experiments with the same seeds reproduce both results,
so they are a property of the estimator (five iterations, the GBDT configuration of
\S\ref{sec:method}), not of the background treatment. This regularization bias is not
included in any quoted uncertainty. On the data, the two background treatments differ by
a median of \SI{0.36}{\percent} and at most \SI{1.4}{\percent} over the reported
functionals. This measures method sensitivity, not a bias. The negative-weight treatment
is not adopted.
% Source: docs/orchestration/OUTCOME-20260925-s5n-stage1-development-fail.md (VL151, VL152);
% KNOWN_ISSUES 75, 77; independently reproduced (state/s5n/stage1/review-round-1.md).

\paragraph{Diagnosis of the departure response; a refinement-capacity candidate.}
An instrumented rerun recorded every OmniFold iteration of these experiments. It
reproduced the development results bitwise and repeated them through the unmodified
driver, which shows the same response to the $\Eavail$ departure (highest-$W$ cell
\SI{73.8}{\percent} against \SI{73.5}{\percent}).

The two steps behave as intended. The first step matches the pseudo-data at reco level
within statistical noise, and the second reproduces its truth-level weight sums to
\SI{0.5}{\percent}. Yet the unfolded result folds back onto the data less well than the
true model does. On noise-free pseudo-data it reproduces more than \SI{99}{\percent} of
the departure's reco-level signal while the $(\Eavail,W)$ cells stay \SIrange{5}{11}{\percent}
off in the median and up to \SI{74}{\percent} in the worst cell. The remaining reco-level
imprint is small relative to the departure but still statistically significant at the
analysis exposure.

Thirty iterations reduce the $\Eavail$ residual by \SI{26}{\percent}, and it is still
falling slowly; the $q_3$ residual grows with iterations. A larger classifier improves the
reco-level fit without improving the $\Eavail$ truth residual, and removing the regression
for events that fail reconstruction does not help it. No single cause is established, and
no authorized estimator change passes its predeclared screen for the physical departure.

The residual background bias has a named cause: the refinement classifier underfits.
Raising its capacity (400 trees, 31 leaves) reduces the background-related bias by more
than \SI{80}{\percent} in every affected cell. A candidate built on this passes the
nominal background-inclusive closure on 40 fresh pseudo-experiments (largest deviation
2.7 standard errors) and a calibration test of its statistical intervals (pull width 1.02).
It fails the predeclared stability test on the data: perturbing the inputs within
floating-point rounding moves the reported functionals by 0.8 statistical standard
deviations in the median and 2.6 at most. The same instability, at most \SI{1.2}{\percent}, explains the
difference between the driver and development pipelines on the data; whether the
data-bootstrap uncertainty already absorbs it is not established. Under withheld
generator shapes the unfolding-model dependence is \SIrange{2}{6}{\percent} in the median
and \SIrange{16}{31}{\percent} at most. Nothing is adopted.
% Source: docs/orchestration/DIAGNOSIS-20260926-s5e-oi192-estimator.md (VL153);
% docs/orchestration/OUTCOME-20260926-s5e-oi192-diagnosis-and-candidate.md (VL154); KNOWN_ISSUES 75, 77, 79, 80;
% independently reproduced (state/s5e/diag/review-round-1.md, state/s5e/cand/review-round-2.md).

\paragraph{Unbinned diagnostic (C2ST).}
As a binning-free check that the unfold removes the data\,$\leftrightarrow$\,MC mismatch, we use
a classifier two-sample test (C2ST): a classifier is trained to separate the data from the
(reweighted) simulation in the full unbinned feature space, and its held-out AUC/accuracy is
compared to the chance value $0.5$. The discriminator AUC drops from $\approx0.535$ before
unfolding to $\approx0.501$ after, indicating that the reweighting removes the reconstructed
data--MC mismatch. We report this strictly as a \emph{descriptive} diagnostic and do \emph{not}
quote a calibrated $p$-value: the C2ST as run trains and evaluates on the same sample without
cross-fitting, and a valid null would require re-running the full (cross-fitted) OmniFold
pipeline per pseudoexperiment. Calibrated unbinned goodness-of-fit remains an
open problem~\cite{unbinnedguide}.

\paragraph{Toy residuals and coverage status.}
The existing 200 2D closure toys fluctuate both the unfolded result and the
stored \texttt{hTruthXSec2D} through the MC bootstrap. The stored truth varies
in all 205 reported bins, so its all-toy mean is not an independent fixed truth
for a split-sample coverage claim. We therefore quote no 2D coverage number
from this ensemble. Its same-ensemble figures
(\SI{\pullFrac}{\percent}, $\langle|r|\rangle=\pullMeanAbs$ vs
$\sqrt{2/\pi}=\pullGauss$, and \SI{97.6}{\percent} of bins above
\SI{65}{\percent}) are retained only as standardized-pull/Gaussianity
diagnostics. A coverage test requires an independently stored unfluctuated
truth reference (or a redesigned ensemble) and remains open.

\paragraph{Stability.}
5-iter total $\sigma$ within \SI{0.026}{\percent} of 10-iter (per-bin shape RMS
\SI{1.54}{\percent}); HistGBT and exact-GBT agree on total $\sigma$ to
\SI{0.04}{\percent}. Efficiency (Fig.~\ref{fig:eff}) reproduces the published
Fig.~5.

\paragraph{Classifier choice (GBDT vs.\ NN).}
Published OmniFold analyses use neural-network
classifiers~\cite{Andreassen:2019cjw,unbinnedguide}; we use gradient-boosted decision
trees, whose output must still be a calibrated posterior for $w=p/(1-p)$ to be valid.
On the step-1 reco classifier (data vs.\ MC reco) the reliability curve is near-diagonal
(Brier \(\approx0.25\) at AUC \(\approx0.537\): data and reco MC are nearly
indistinguishable, so the reweight is a modest correction). A full scalar
Keras-MLP OmniFold cross-check on the same weighted 3D inputs and five-iteration
loop gives total ratio $1.0078$ versus the GBDT reference and median projection
differences of \SIrange{0.7}{1.4}{\percent} (Fig.~\ref{fig:calibration}).
The agreement in the populated bulk establishes that the scalar result is not
specific to the production classifier family. A stronger family cross-check, on a
low-level recoil point-cloud representation rather than the engineered scalars, is
reported with the rest of that track in \S\ref{sec:pet-shape}.

\begin{figure}[H]
  \centering
  \includegraphics[width=0.95\textwidth]{nn_vs_gbdt_full}
  \caption{Full five-iteration scalar OmniFold cross-check using identical
  weighted 3D inputs and extraction, with only the classifier family changed.
  Top: absolutely normalized GBDT and Keras-MLP projections. Bottom: their
  binwise ratio. The MLP/GBDT total-cross-section ratio is $1.0078$; median
  absolute projection differences are \SI{1.20}{\percent} in $\pT$,
  \SI{1.36}{\percent} in $\pPar$, and \SI{0.66}{\percent} in $\Eavail$.
  The larger excursions are confined to the sparsely populated outermost
  $\pT$ and $\pPar$ bins.}
  \label{fig:calibration}
\end{figure}




\begin{figure}[H]
  \centering
  \includegraphics[width=0.62\textwidth]{MEFHC_5iter_xsec_eff_heatmap}
  \caption{Selection efficiency over the $(\pT,\pPar)$ truth grid. Efficiency is
  diagnostic only --- OmniFold's miss treatment performs the acceptance
  correction natively (\S\ref{sec:method-algo}).}
  \label{fig:eff}
\end{figure}
